% Independently audited replacement for Section 10.
\section{Mean-field stability with a prescribed free-sea reference}
\label{sec:meanfield}

For the fully free-Wick-ordered, photon-free Hamiltonian specified earlier,
one can prove a vacuum-relative bound $-CK$ on quasi-free states, including
pair coherence and point nuclei.  The volume range is restricted by an
explicit lower bound on $mL$ as well as the ultraviolet scale.  The result
also holds on exactly neutral Slater determinants.  Mixed quasi-free states
with the prescribed mean charge form a larger variational class; they must
not be described as states in a single exact-charge sector.

\subsection{Quasi-free states, charge, and the energy functional}

In the finite sea representation let $\gamma$ be a one-particle density
matrix with $0\le\gamma\le1$, and write
\[
 Q=\gamma-P_-,\qquad
 \rho_Q(x)=-\tr_{\C^4}Q(x,x),\qquad
 \mathcal G_Z=\{0\le\gamma\le1:\Tr Q=Z_{\rm tot}\}.
\]
It determines a unique gauge-invariant quasi-free state $\omega_\gamma$.
Gauge invariance here means invariance under the global charge symmetry;
the state may have nonzero electron--positron coherence $Q_{+-}$ in the
free particle--hole representation.  It satisfies
\[
 \omega_\gamma(N_e-N_p)=\Tr Q,
 \qquad
 \operatorname{Var}_{\omega_\gamma}(N_e-N_p)
 =\Tr\bigl(\gamma(1-\gamma)\bigr).
\]
Consequently $\gamma\in\mathcal G_Z$ imposes neutrality in expectation.
It imposes exact neutrality precisely when $\gamma$ is a projection of
rank $2M_L+Z_{\rm tot}$, in which case the state is a Slater determinant.
Both classes obey $Z_{\rm tot}\le2M_L$.  No assertion below identifies the
infimum over the relaxed class with the infimum over these projections.

The elementary inequality $\gamma^2\le\gamma$ gives
$Q^2\le Q_{++}-Q_{--}$, and hence
\begin{equation}
 \Tr(D_0Q)
 =\Tr\bigl(|D_0|(Q_{++}-Q_{--})\bigr)
 \ge\Tr(|D_0|Q^2)\ge m\Tr Q^2.
 \label{eq:kin-mf}
\end{equation}
Also $-P_-\le Q\le P_+$ and, by evaluation on band-limited delta vectors,
$|\rho_Q(x)|\le2d_L$.

\begin{lemma}[Hartree--Fock energy for the specified Wick convention]
\label{lem:hf-energy}
For $H^{\rm C}_{\Lambda,L}=H_D+H_C$, with the full free-vacuum Wick ordering
used in the main model, every $0\le\gamma\le1$ satisfies
\begin{multline}
 \omega_\gamma(H^{\rm C}_{\Lambda,L}(R,Z))
 =\mathcal E_{R,Z}(\gamma)
 :=\Tr(D_0Q)-\frac\alpha2\mathrm{Ex}(Q)
       +\frac\alpha2D_L(\rho_Q,\rho_Q)\\
       +\alpha D_L(\rho_Q,\nu)
       +\alpha\sum_{i<j}Z_iZ_jG_L(R_i-R_j),
 \label{eq:hf-energy}
\end{multline}
where
\[
 \mathrm{Ex}(Q)=\iint G_L(x-y)\|Q(x,y)\|_{\HS}^2\dd x\dd y.
\]
This exchange expression uses the full zero-mean periodic Green function;
it is not asserted to be nonnegative.
\end{lemma}
\begin{proof}
The kinetic expectation is $\Tr(D_0Q)$.  The quasi-free Wick identity
\eqref{eq:trial-wick} is valid also for mixed quasi-free states.  Its direct
term is $D_L(\rho_Q,\rho_Q)$, whereas its exchange term is
\[
 \frac1V\sum_{k\ne0}\frac{4\pi}{|k|^2}\Tr(QT_{-k}QT_k)
 =\iint G_L(x-y)\tr\bigl(Q(x,y)Q(y,x)\bigr)\dd x\dd y.
\]
Since $Q(y,x)=Q(x,y)^*$, this is $\mathrm{Ex}(Q)$.  Only finitely many
Fourier coefficients contribute.  The linear nuclear coupling has the
stated expectation.  This proves the formula and fixes its signs.
\end{proof}

The disappearance of a linear exchange self-energy in
\eqref{eq:hf-energy} depends on full Wick ordering relative to $P_-$.
It would not hold unchanged for a Hamiltonian with no normal ordering, or
for density-product ordering with only a scalar vacuum subtraction.

\subsection{A periodic Coulomb--Kato inequality with explicit constants}

Periodic Coulomb--Kato estimates are established tools: see, in particular,
Lemma~1 of Hainzl--Lewin--Solovej~\cite{HLSol2007}.  The following proof
records an explicit constant for the zero-mean Green function used here.
Their periodic potential is $W_L=G_L+\mu/L$, normalized to have minimum
zero; their lemma gives a bounded $O(L^{-1})$ comparison with $1/|x|$ and
a cutoff Coulomb--Dirac form constant tending to its infinite-volume
value, which is at most $\pi/2$.  The additive constant below is an
explicit adaptation, not a quoted constant from that result.

\begin{lemma}[Periodic Coulomb--Kato inequality]
\label{lem:pkato}
Set
\[
 c_*:=6+\frac{3\pi\sqrt6}{2}<18.
\]
For every $L>0$, $y\in\T_L^3$, and
$f\in H^{1/2}(\T_L^3;\C^4)$,
\begin{equation}
 \int_{\T_L^3}G_L(x-y)|f(x)|^2\dd x
 \le\frac\pi2\langle f,|p|f\rangle
       +\frac{c_*}{L}\|f\|^2.
 \label{eq:pkato}
\end{equation}
\end{lemma}
\begin{proof}
Translation reduces to $y=0$.  We first justify, without numerical evidence,
the convenient comparison
\begin{equation}
 G_1(x)-|x|^{-1}<-\frac1{20}
 \quad\text{on }[-1/2,1/2]^3,
 \label{eq:ewald-certified}
\end{equation}
with the value at zero interpreted by continuity.  The heat representation
of $G_1$, split at $t=1/36$, and Poisson summation give the Ewald identity
\begin{equation}
 G_1(x)=\sum_{n\in\Z^3}\frac{\operatorname{erfc}(3|x+n|)}{|x+n|}
 +\frac1\pi\sum_{k\ne0}
      \frac{e^{-\pi^2|k|^2/9}}{|k|^2}\cos(2\pi k\cdot x)
 -\frac\pi9.
 \label{eq:ewald-audit}
\end{equation}
Both sums converge absolutely away from the isolated displayed singularity.
Here $\operatorname{erfc}(u)=\frac2{\sqrt\pi}\int_u^\infty e^{-v^2}\dd v$.
The elementary bound
$\operatorname{erfc}(u)\le e^{-u^2}/(\sqrt\pi u)$ gives, uniformly on the
cell, the following bound for the nonzero real-space images:
\begin{align}
 I&:=\sum_{n\ne0}\frac{\operatorname{erfc}(3|x+n|)}{|x+n|}\nonumber\\
 &\le\frac8{\sqrt\pi}(e^{-9/4}+e^{-9/2})
       +\frac{32}{9\sqrt\pi}e^{-27/4}
       +\frac{9e^{-81/4}}{1-e^{-36}}
 <\frac8{15}.
 \label{eq:ewald-real-bound}
\end{align}
Indeed the shell $|n|_\infty=1$ has $6,12,8$ points with respectively one,
two, and three nonzero coordinates, whose distances from the cell are at
least $1/2,\sqrt2/2,\sqrt3/2$.  For $|n|_\infty=j\ge2$ there are
$24j^2+2$ points and $|x+n|\ge j-1/2$; their contribution is at most
$9e^{-9(j-1/2)^2}$.  Successive exponents differ by at least $36$.
For a fully rational check of the last strict inequality, use
\[
 z=\frac{53}{500},\qquad e^{-9/4}<z,\qquad\sqrt\pi>\frac{177}{100},
\]
and observe that
\[
 \frac{800}{177}(z+z^2)+\frac{3200}{1593}z^3
       +\frac{9z^9}{1-z^{16}}<\frac8{15}.
\]
The exponential bound follows from
$\sum_{j=0}^7(9/4)^j/j!>500/53$.

For the absolute value of the reciprocal-space sum, set $t=669/2000$.
The elementary bounds $\pi>157/50$ and
\[
 \sum_{j=0}^8\frac{(24649/22500)^j}{j!}>\frac{2000}{669}
\]
imply $e^{-\pi^2/9}<t$.  The multiplicities of $|k|^2$ in
$[-2,2]^3\cap\Z^3\setminus\{0\}$ are
\[
 (s,N_s)=(1,6),(2,12),(3,8),(4,6),(5,24),
          (6,24),(8,12),(9,24),(12,8).
\]
The remaining shell $|k|_\infty=j\ge3$ contributes at most $26t^{j^2}$
before the factor $1/\pi$.  Thus
\begin{align}
 F&:=\frac1\pi\sum_{k\ne0}\frac{e^{-\pi^2|k|^2/9}}{|k|^2}\nonumber\\
 &<\frac{50}{157}\left(
 6t+6t^2+\frac83t^3+\frac32t^4+\frac{24}{5}t^5+4t^6
 +\frac32t^8+\frac83t^9+\frac23t^{12}
 +\frac{26t^9}{1-t^7}\right)<\frac9{10}.
 \label{eq:ewald-fourier-bound}
\end{align}
All displayed decimal-free comparisons are inequalities between rational
numbers or follow from the stated finite Taylor sums.

The function $r\mapsto\operatorname{erf}(3r)/r$ is decreasing.  Since
$r\le\sqrt3/2$, the same complementary-error-function bound gives
\[
 \frac{\operatorname{erf}(3r)}r
 \ge\frac2{\sqrt3}\operatorname{erf}(3\sqrt3/2)
 >\frac{15}{13}\frac{2999}{3000}>\frac{23}{20}.
\]
Here we used $2/\sqrt3>15/13$,
$3\sqrt{3\pi}/2>4$, and $z^3/4<1/3000$.
Combining this with $\pi/9>1/3$,
\eqref{eq:ewald-real-bound}, and \eqref{eq:ewald-fourier-bound} proves
\eqref{eq:ewald-certified}, since
$8/15+9/10-1/3-23/20=-1/20$.  In particular
$G_L(x)\le1/|x|$ on the centered cell after scaling.

We next give the localization argument, including the folding step.  Define
the nonnegative Lipschitz function
\[
 g(t)=\begin{cases}
 1,&|t|\le1/6,\\
 \cos\bigl(\frac{3\pi}4(|t|-1/6)\bigr),&1/6<|t|<5/6,\\
 0,&|t|\ge5/6,
 \end{cases}
\]
and set $\phi(x)=\prod_{i=1}^3g(x_i/L)$,
$\phi_n(x)=\phi(x-Ln)$.  The cosine--sine identity on each overlap gives
$\sum_n\phi_n(x)^2=1$.  Lipschitz regularity is sufficient for multiplication
on $H^{1/2}$ and for the integrable remainder below; smoothness is not needed.
With $\mathcal C=[-L/2,L/2]^3$ and periodic extension of $f$, every point of
$\mathcal C\cap\operatorname{supp}\phi_n$ with $n\ne0$ has $|x|\ge L/6$.
Consequently the Euclidean Coulomb--Kato form inequality
$|x|^{-1}\le(\pi/2)|p|$ (see \cite{Herbst1977}) gives
\begin{equation}
 \int_{\mathcal C}\frac{|f(x)|^2}{|x|}\dd x
 \le\frac\pi2\langle\phi f,|p|_{\R^3}\phi f\rangle
       +\frac6L\|f\|^2.
 \label{eq:pkato-local}
\end{equation}

With the normalization of $|p|$ used here its torus quadratic form is
\begin{equation}
 \langle f,|p|f\rangle_{\T_L^3}
 =\frac1{2\pi^2}\int_{\mathcal C}\dd x\int_{\R^3}\dd y\,
            \frac{|f(x)-f(y)|^2}{|x-y|^4}.
 \label{eq:torus-half-laplacian}
\end{equation}
The coefficient is checked on Fourier modes, or by periodizing the
Euclidean quadratic form.  Folding the $x$ integral in the Euclidean form
for $\phi f$ over cells, while translating $y$ by the same lattice vector,
gives the exact identity
\[
 \langle\phi f,|p|_{\R^3}\phi f\rangle
 =\langle f,|p|f\rangle_{\T_L^3}+R,
\]
where
\[
 R=\frac1{2\pi^2}\int_{\mathcal C}\dd x\int_{\R^3}\dd y\,
 \frac{\sum_n(\phi_n(x)-\phi_n(y))^2
             \operatorname{Re}\langle f(x),f(y)\rangle}{|x-y|^4}.
\]
This follows directly by expanding the squares and using
$\sum_n\phi_n^2=1$.  The vector $\Phi(x)=(\phi_n(x))_n$ is a unit vector
with nonnegative entries and
\[
 \sum_n|\nabla\phi_n(x)|^2\le\frac{3(3\pi/4)^2}{L^2}.
\]
Thus, with $a=3\pi\sqrt3/4$,
\[
 \sum_n(\phi_n(x)-\phi_n(y))^2
 \le\min\{a^2|x-y|^2/L^2,2\}.
\]
Use $2|f(x)||f(y)|\le|f(x)|^2+|f(y)|^2$.  The two resulting integrals
are equal after exchanging variables and folding: the kernel is symmetric
and invariant under simultaneous lattice translations.  Therefore
\begin{align*}
 |R|&\le\frac{\|f\|^2}{2\pi^2}
       \int_{\R^3}\frac{\min\{a^2|z|^2/L^2,2\}}{|z|^4}\dd z\\
 &=\frac{4\sqrt2a}{\pi L}\|f\|^2
  =\frac{3\sqrt6}{L}\|f\|^2.
\end{align*}
The identities can first be proved for smooth periodic $f$ and then extended
by density.  Combining this estimate, \eqref{eq:pkato-local}, and
$G_L\le1/|x|$ proves \eqref{eq:pkato} with the stated $c_*$.
\end{proof}

\begin{corollary}[Exchange bound]
\label{cor:exchange}
For every self-adjoint $Q$ on the cutoff one-particle space,
\begin{equation}
 \mathrm{Ex}(Q)
 \le\frac\pi2\Tr(|p|Q^2)+\frac{c_*}{L}\Tr Q^2.
 \label{eq:exchange}
\end{equation}
For $Q=\gamma-P_-$ with $0\le\gamma\le1$, the right side is at most
$\bigl(\pi/2+c_*/(mL)\bigr)\Tr(D_0Q)$.
\end{corollary}
\begin{proof}
Apply \eqref{eq:pkato} to $f(x)=Q(x,y)e_b$, sum over the spin index $b$,
and integrate over $y$.  The two right-side expressions become
$\Tr(|p|Q^2)$ and $\Tr Q^2$.  Then use $|p|\le|D_0|$ and
\eqref{eq:kin-mf}.
\end{proof}

\subsection{The mean-field theorem and its volume range}

\begin{theorem}[Mean-field relative stability with pairs and point nuclei]
\label{thm:mf}
Let $0\le\alpha<4/\pi$ and
\[
 \theta_L=1-\frac{\alpha\pi}4-\frac{\alpha c_*}{2mL},
 \qquad c_*=6+\frac{3\pi\sqrt6}{2}.
\]
Assume $L\ge\Lambda_f^{-1}$ and
\begin{equation}
 L>\frac{\alpha c_*}{2m(1-\alpha\pi/4)},
 \quad\text{equivalently }\theta_L>0.
 \label{eq:mf-volume-range}
\end{equation}
For distinct point nuclei with integral $1\le Z_i\le Z_{\max}$ and nonempty
$\mathcal G_Z$, every $\gamma\in\mathcal G_Z$ satisfies
\begin{equation}
 \mathcal E_{R,Z}(\gamma)
 \ge\theta_L\Tr(D_0Q)-A^{\rm MF}_\Lambda K,
 \qquad
 A^{\rm MF}_\Lambda=C\alpha\Lambda_f(Z_{\max}+Z_{\max}^2),
 \label{eq:mf-bound}
\end{equation}
with a universal finite constant $C$.  Without nuclei,
\[
 \mathcal E_0(\gamma)
 \ge\theta_L\Tr(D_0Q)+\frac\alpha2D_L(\rho_Q,\rho_Q)\ge0,
\]
with equality only at $\gamma=P_-$.  Thus the source-free Hartree--Fock
vacuum for this Wick-ordered model is exactly $P_-$ and has energy zero.
Both the relaxed mean-charge infimum and the infimum over exactly neutral
Slater determinants obey
\begin{equation}
 E^{\rm HF}(R,Z)-E^{\rm HF,\vac}_{\Lambda,L}
 \ge-A^{\rm MF}_\Lambda K.
 \label{eq:mf-relative}
\end{equation}
The constant in this relative bound is independent of $L$ within
\eqref{eq:mf-volume-range}, of the nuclear positions, and of the pair
coherence.  The assertion does not extend to all $L>0$.
\end{theorem}
\begin{proof}
The exchange bound gives
$\Tr(D_0Q)-\frac\alpha2\mathrm{Ex}(Q)\ge\theta_L\Tr(D_0Q)$.
The remaining three terms of \eqref{eq:hf-energy} are the classical
density-product electrostatic expression.  Apply the heat decomposition
\eqref{eq:heat-identity}, now to the scalar density $\rho_Q$.
Its square is nonnegative, $|\rho_Q|\le2d_L$, and
$Z_{\rm tot}\le2d_LV$.  Exactly as in the operator estimate, this gives
\[
 \frac\alpha2D_L(\rho_Q,\rho_Q)+\alpha D_L(\rho_Q,\nu)
   +\alpha\sum_{i<j}Z_iZ_jG_L(R_i-R_j)
 \ge-C\alpha\bigl(d_Ls Z_{\max}+s^{-1/2}Z_{\max}^2\bigr)K.
\]
Set $s=\Lambda_f^{-2}$ and use $d_L\le C\Lambda_f^3$ for
$L\ge\Lambda_f^{-1}$.  This proves the estimate on the relaxed class,
and therefore also on its subset of neutral projections.
For $K=0$ retain the positive direct term.  Vanishing of the energy forces
$\Tr(D_0Q)=0$, hence $Q=0$ by \eqref{eq:kin-mf}.  The free sea is itself
an admissible neutral projection, so both vacuum infima equal zero.
\end{proof}

\begin{remark}[Scope, ordering, and prior work]
\label{rem:mf-scope}
\begin{enumerate}[label=(\alph*)]
\item The volume term in the general many-body Coulomb lower bound came
from bounding Wick-reordering corrections in operator norm.  On quasi-free
states their exact expectation combines into the exchange functional,
which is controlled by the kinetic energy.  This is why that volume term
can be removed in the present variational class.
\item For the physical coupling, \eqref{eq:mf-volume-range} is an explicit
microscopic lower bound on $L$; it is not an assumption that can be omitted
from the theorem.  The restriction $L\ge\Lambda_f^{-1}$ is independent of
it.  The relative bound remains uniform as $L\to\infty$ at fixed cutoff.
\item No restriction on $Z_{\max}\alpha$ is needed for this fixed-cutoff
argument.  The price is $A^{\rm MF}_\Lambda=O(\alpha\Lambda_fZ_{\max}^2)$.
A cutoff-uniform point-nucleus assertion requires new estimates and may
require a subcritical charge condition.  The present proof provides none.
\item The statement that $P_-$ is the Hartree--Fock vacuum is specific to
the Hamiltonian fully Wick ordered relative to that projector.  The
un-normal-ordered model of Hainzl--Lewin--Solovej~\cite{HLSol2007} instead
constructs a self-consistent interacting vacuum and its associated BDF
functional through a thermodynamic limit.  Their periodic Kato estimate
is prior art, and their vacuum is not to be identified with $P_-$ by
changing notation.  The stable polarized-vacuum results
\cite{HLS2005,HLS2008} likewise require their stated reference operators
and external-density assumptions.
\item Bach--Barbaroux--Helffer--Siedentop~\cite{BBHS1998} already prove
stability for a Hartree--Fock electron--positron model with point nuclei
and classical magnetic fields, including physical charges up to $67$.
Their reference splitting uses a Dirac operator with a configuration-dependent
effective nuclear potential (and, in the magnetic theorem, the classical
vector potential); it is not the fixed free-Wick-ordered torus Hamiltonian above.
This section supplies an explicit estimate for the chosen cutoff and
normalization, rather than a claim that stability with pairs and point
nuclei was previously unknown.
\item A lower bound on quasi-free expectations does not bound arbitrary
correlated states: the variational inequality goes in the opposite
direction. The correlated-state theorems above use an independent
canonical comparison and their own effective-mass hypothesis.
\end{enumerate}
\end{remark}

\subsection{Mean field with transverse photons}
\label{subsec:mf-photons}

Restrict to product states $\omega_\gamma\otimes\Phi$ in the
positive-photon-mass model \eqref{eq:Hct}; entangled fermion--photon states
are not included.  After the photon squeezing transformation of
Lemma~\ref{lem:proca}, minimization over photon states is exact and yields
\begin{equation}
 \inf_\Phi(\omega_\gamma\otimes\Phi)(H^{(\delta\mu)})
 =\mathcal E_{R,Z}(\gamma)-\mathcal M_{\delta\mu}(Q)
       +E_\gamma(\delta\mu),
 \label{eq:mf-photon}
\end{equation}
where
\[
 \mathcal M_{\delta\mu}(Q)
 =\frac{e^2}{2V}\sum_{k,\lambda}
       \frac{|\varepsilon_\lambda(k)\cdot\widehat{j_Q}(k)|^2}
            {|k|^2+\delta\mu^2},
 \qquad
 j_Q(x)=\tr_{\C^4}(\boldsymbol{\mathsf a}Q(x,x)).
\]
The minimizing photon state is coherent in the squeezed representation;
the original representation includes the corresponding fixed squeezing.

\begin{proposition}[Mean-field relative stability with a sufficiently large bare photon mass]
\label{prop:mf-heavy}
Assume Theorem~\ref{thm:mf} and
\[
 \eta_L:=\frac{2e^2d_L}{m\delta\mu^2}<\theta_L.
\]
Every product state with $\gamma\in\mathcal G_Z$ then satisfies
\[
 (\omega_\gamma\otimes\Phi)(H^{(\delta\mu)})-E_\gamma(\delta\mu)
 \ge(\theta_L-\eta_L)\Tr(D_0Q)-A^{\rm MF}_\Lambda K.
\]
The product-state vacuum energy is $E_\gamma(\delta\mu)$, attained at
$P_-\otimes U^*\Omega_\gamma$.  This gives the relative bound $-A^{\rm MF}_\Lambda K$
both in the relaxed mean-charge class and on exactly neutral Slaters
tensored with photon states.
\end{proposition}
\begin{proof}
The matrices $\mathsf a_i/2$ are orthonormal in the Hilbert--Schmidt inner
product.  Bessel's inequality therefore gives, for each $4\times4$ matrix $M$,
\[
 \sum_i|\tr(\mathsf a_iM)|^2\le4\|M\|_{\HS}^2.
\]
At fixed $x$, the four band-limited spin delta vectors are mutually
orthogonal with squared norm $d_L$.  A second application of Bessel gives
\[
 \sum_{a,b}|\langle\delta_x^a,Q\delta_x^b\rangle|^2
 \le d_L\sum_b\|Q\delta_x^b\|^2.
\]
Integration uses the resolution of the identity by these delta vectors and
proves
\[
 \int|j_Q(x)|^2\dd x
 \le4\int\|Q(x,x)\|_{\HS}^2\dd x
 \le4d_L\Tr Q^2
 \le\frac{4d_L}{m}\Tr(D_0Q).
\]
Consequently
$\mathcal M_{\delta\mu}(Q)\le
e^2(2\delta\mu^2)^{-1}\int|j_Q|^2\le\eta_L\Tr(D_0Q)$.
Apply Theorem~\ref{thm:mf} and use the free sea as the vacuum trial state.
The estimate uses both spin and spatial Bessel inequalities.
\end{proof}

This is a coercivity statement for a positive bare $A^2$ coefficient in the
modified model.  It does not determine a nonperturbative photon pole mass
or prove a physical mass-renormalization condition.

\begin{proposition}[Failure of strict kinetic absorption at the statically matched mass]
\label{prop:ward-absorption-false}
Fix $\alpha>0$, $m>0$, and positive cutoffs, and set
$\delta\mu^2=\mu_{\rm ind}^2=2e^2\chi(0)$.
There is no constant $\vartheta<1$, uniform for large $L$, such that
\begin{equation}
 \mathcal M_{\mu_{\rm ind}}(Q)
 \le\vartheta\Tr(D_0Q)
 \quad\text{for all exactly neutral Slater projectors }\gamma.
 \label{eq:ward-absorption-false}
\end{equation}
In particular the proposed sufficient estimate with a coefficient
$\vartheta_L<\theta_L\le1-\alpha\pi/4$ and an error $CK$ is false already
at $K=0$.
\end{proposition}
\begin{proof}
Take $L$ large enough that $2\pi/L\le\Lambda_\gamma$, so the photon mode
used below is retained. Choose $q=2\pi e_1/L$, $\varepsilon=e_2$, and
$A=P_f(\varepsilon\cdot\boldsymbol{\mathsf a})\cos(q\cdot x)P_f$.
Let $X=P_+AP_-$ and define $B:P_-\mathfrak h_f\to P_+\mathfrak h_f$ in
free eigenbases by
\[
 B_{ij}=\frac{X_{ij}}{E_i+|E_j|}.
\]
Regard $S=B-B^*$ as an anti-self-adjoint operator on the full one-particle
space, and put
\[
 \gamma_t=e^{tS}P_-e^{-tS},\qquad Q_t=\gamma_t-P_-.
\]
These are exactly neutral Slater projectors for every real $t$.
With
\[
 s_L=\sum_{i\in+,j\in-}\frac{|X_{ij}|^2}{E_i+|E_j|}
     =\frac V2\chi_L(q;\varepsilon),
\]
the elementary expansion of the unitary rotation gives
\[
 \Tr(D_0Q_t)=t^2s_L+o(t^2),\qquad
 J_t:=\Tr(AQ_t)=2ts_L+O(t^2).
\]
The first identity follows because the second-order diagonal blocks of
$Q_t$ are $t^2BB^*$ and $-t^2B^*B$; the first-order off-diagonal block is
$tB$.  Optimizing just the single real transverse field
$A_{\rm cl}=-a\varepsilon\cos(q\cdot x)$ gives the lower bound
\[
 \mathcal M_{\delta\mu}(Q_t)
 \ge\frac{e^2J_t^2}{V(|q|^2+\delta\mu^2)}.
\]
It follows that
\begin{equation}
 \liminf_{t\to0}
 \frac{\mathcal M_{\delta\mu}(Q_t)}{\Tr(D_0Q_t)}
 \ge\frac{2e^2\chi_L(q;\varepsilon)}{|q|^2+\delta\mu^2}.
 \label{eq:ward-ratio}
\end{equation}
By Proposition~\ref{prop:response}, the right side tends to $1$ as
$L\to\infty$ at the statically matched value.  For any fixed $\vartheta<1$, first
choose $L$ so that this right side exceeds $\vartheta$, and then choose a
sufficiently small nonzero $t$.  This contradicts
\eqref{eq:ward-absorption-false}.  Since
$\theta_L\le1-\alpha\pi/4<1$, the last assertion follows as well.
\end{proof}

\begin{remark}[The actual statically matched mean-field question]
\label{rem:mf-photon-open}
The preceding proposition converts the previously proposed strict
absorption inequality from an open question to a disproved statement.
It does not decide stability of the complete mean-field functional.
The direct Coulomb term, exchange, and the correct interacting mean-field
vacuum have to be retained in such an argument.

At fixed cutoff the sufficient large-mass condition has continuum-volume
ratio $\eta_L\to E_\Lambda/m$ at the static matching value, so it does not establish
that case.  More importantly, Proposition~\ref{prop:ward-absorption-false}
shows that merely improving it to a strict uniform kinetic absorption
below one cannot close the argument.

For a precise remaining question define, using the same charge class on
both sides,
\[
 F_{R,Z}^{\delta\mu}(\gamma)
 =\mathcal E_{R,Z}(\gamma)-\mathcal M_{\delta\mu}(Q),
 \qquad
 f_0^{\delta\mu}=\inf_{\gamma\in\mathcal G_0}F_0^{\delta\mu}(\gamma).
\]
The desired statically matched product-state statement is
\[
 \inf_{\gamma\in\mathcal G_Z}F_{R,Z}^{\mu_{\rm ind}}(\gamma)
       -f_0^{\mu_{\rm ind}}\ge-C_\Lambda K,
\]
or its separate version with both infima restricted to the appropriate
exact-charge projections.  Neither version is proved here.  In particular
$f_0^{\mu_{\rm ind}}$ must not be set to zero using the large-mass theorem,
whose hypothesis fails at this tuning.  These remain mean-field questions
for a modified photon-mass model, not proofs of the original massless
many-body QED assertion.
\end{remark}
